\section{Partial traces of regrouped tensor products}
\label{sec:tensor_power_partial_trace}

\begin{definition}[Collecting the factors of each copy]
\label{def:copies_product_equiv}
\lean{TensorPower.copiesProductEquiv}
Given a bipartite coordinate decomposition $A\simeq Q\times T$, apply
this decomposition in each of $m$ copies and collect the two families:
\begin{align*}
 A^m\simeq Q^m\times T^m.
\end{align*}
The order of the copies is the same in both families.
\end{definition}

\begin{theorem}[Partial trace of a tensor product]
\label{thm:partial_trace_tensor_power_reindex}
\lean{Matrix.partialTraceRight_finKronecker_reindex}
For matrices $M_1,\ldots,M_m$ on $A$, their tensor product in the
preceding coordinates satisfies
\begin{align*}
 \operatorname{Tr}_{T^m}\Bigl(\bigotimes_{j=1}^m M_j\Bigr)
 =\bigotimes_{j=1}^m\operatorname{Tr}_T M_j.
\end{align*}
The formula includes $m=0$ and requires no positivity or normalization.
It is the tensor-product identity used in
\cite[Section~7, lines~550--560]{OpenAI2026AreaLaw}.
\end{theorem}
\begin{proof}
\uses{def:copies_product_equiv}
Expand one matrix entry of the partial trace. The sum over all
complementary configurations is the product of the individual sums.
\end{proof}
